\section{To Die in Christ}

We shall die, and the hour of death will end the time given us for repentance. The Church teaches us to keep that hour before us and to prepare for it in Christ. Through Baptism we have already died with him sacramentally; to die in his grace is to complete that dying with him and enter his life (CCC 1010--1014). We seek a good death by living faithfully in that grace and asking Christ to preserve us in it to the end.

Someone else can stand beside our deathbed, but no one can undergo our dying for us. A priest can absolve and anoint; friends can forgive and keep watch; the Church can place her prayer around our weakness. These are real gifts, and our need for them can be very great. Yet the life being surrendered is ours. The accumulated provisions by which we have made ourselves secure will cease to be available in their familiar way. Possessions remain behind. Offices pass to successors. Unfinished plans fall into other hands. Even the body through which we have known the world will be taken from our keeping.

We can begin by considering the good we have postponed, the sin we have excused, and the forgiveness we have withheld. There is still time to seek pardon and amend our conduct. We need not wait until death becomes imminent before examining how we are living. The same Lord who will receive our last breath is calling us now, through his word, the sacraments, and the duties nearest at hand.

The Four Last Things---death, judgment, hell, and heaven---concern the eternal destiny of this life. Death leads to judgment; heaven and hell are the opposed definitive outcomes. Purgatory is the final purification of those already assured of heaven. The person who acts today will stand before God, with the good received from his grace and the sins for which mercy is needed. Our daily choices matter because we are choosing how to respond to the Lord with whom we are called to live forever.

The Church's year gives this contemplation a familiar setting. At All Saints we honor those who have attained holiness; at All Souls we pray for the departed who still need our prayers. In Advent we await Christ's coming, the resurrection of the dead, and the completion of history. Another year of our life has been spent when these seasons return. Someone who once prayed beside us may now be among those for whom we pray. A burden once remote may have become our daily responsibility. Remembering death again allows us to consider these changes before God and examine how we have responded to them.

We cannot command another year, another day, or even another opportunity to do what conscience already requires. Christ has died and risen, and he offers us the grace to turn to him now. These truths give us reason both to repent promptly and to hope. We can spend much time imagining our last hours while avoiding that present response. Prayer about death bears fruit when it leads us to seek the Lord's mercy and follow him in the life still given to us.

We can prepare without knowing what our death will look like. It may come after a long illness or with little warning; it may leave time for speech or take away that power. We cannot arrange the entire event as a final expression of our character. We can ask for perseverance, seek reconciliation, learn charity, and consent today to belong to Christ in circumstances we cannot yet foresee. When we cannot imagine how we would endure them, we can ask for the help he gives to the weak. The Church teaches us to entrust the unknown hour to him and fulfill the duties already entrusted to us.

